% SGA 3, Expose VII_B, English reconstruction -- Loop 1.
% Scope: Section 5.1.5 in full on authority local pages 88--89;
% stops immediately before Section 5.2.
% Combined-reader pages: 601--602. Re-edition printed pages: 590--591.
% Editorial note in scope: 152.
% Authority: Polo--Gille Exp7B-13oct24.pdf,
% SHA-256 D303690E0195654BBE7C3DAA5D89EA23E24CD709A3575217234AC19AE8836694.
% Jacob Reinhold's CC BY 4.0 Markdown was consulted as comparison lineage.
% No genuine diagram occurs in this component; all mathematics is native LaTeX.

\subsubsection*{5.1.5}
\phantomsection
\label{sga3:VIIB:5.1.5}

\medskip
\noindent\textbf{引理 5.1.5.}
\phantomsection
\label{sga3:VIIB:5.1.5:filtered-groups-lemma}
\footnote{\label{sga3:VIIB:nde:152}原文只陈述了 \textup{\SGARefLink{sga3:VIIB:5.1.5:item:ii}{(ii)}}。为方便读者，编者在 \textup{\SGARefLink{sga3:VIIB:5.1.5:item:i}{(i)}} 中陈述了 \textup{[CA]}，\S~V.7 的引理~1。}
{\itshape
\begin{enumerate}
\item[\textup{\SGAInlineRefTarget{sga3:VIIB:5.1.5:item:i}(i)}]
设 \(M\) 和 \(N\) 为两个阿贝尔群，分别赋以由子群递减序列
\[
  (M_n)_{n\in\mathbb Z},
  \qquad
  (N_n)_{n\in\mathbb Z}.
\]
给出的滤过。假设
\[
  \bigcup_{n\in\mathbb Z}M_n=M,
  \qquad
  \bigcup_{n\in\mathbb Z}N_n=N,
\]
及
\[
  \bigcap_{n\in\mathbb Z}M_n=0,
  \qquad
  \bigcap_{n\in\mathbb Z}N_n=0,
\]
且 \(M\) 关于 \(M_n\) 所定义的拓扑完备。设
\[
  f:M\longrightarrow N
\]
为滤过群同态。

\begin{enumerate}
\item[\textup{a)}]
若 \(f\) 在相联系的分次群上诱导满射，则 \(f\) 为满射，且 \(N\) 关于 \(N_n\) 所定义的拓扑完备。

\item[\textup{b)}]
若 \(f\) 在相联系的分次群上诱导单射，则 \(f\) 为单射。
\end{enumerate}

\item[\textup{\SGAInlineRefTarget{sga3:VIIB:5.1.5:item:ii}(ii)}]
设 \(B\) 为阿贝尔群，
\[
  B=B'_0\supset B'_1\supset\cdots,
  \qquad
  B=B_0\supset B_1\supset\cdots
\]
为 \(B\) 的两个由子群给出的分离滤过，满足
\[
  B'_n\subset B_n
  \qquad\text{对每个 }n.
\]
假设 \(B\) 关于滤过 \((B'_n)\) 所定义的拓扑完备。若映射
\[
  B'_i/B'_{i+1}
  \longrightarrow
  B_i/B_{i+1}
\]
对每个 \(i\) 都为满射，则
\[
  B'_n=B_n
  \qquad\text{对每个 }n.
\]
\end{enumerate}
}

\noindent\textit{证明。}
\textup{\SGARefLink{sga3:VIIB:5.1.5:item:i}{(i)}} 即 \textup{[CA]}，\S~V.7 的引理~1\textup{（另见 [BAC]，III，\S~\SGARefLink{sga3:VIIB:2.8}{2.8}）}。取
\[
  M=B'_n\supset B'_{n+1}\supset\cdots,
  \qquad
  N=B_n\supset B_{n+1}\supset\cdots.
\]
即得 \textup{\SGARefLink{sga3:VIIB:5.1.5:item:ii}{(ii)}}。
